\documentclass[10pt]{article}
\usepackage[letterpaper,margin=0.72in,headheight=15pt]{geometry}
\usepackage[T1]{fontenc}
\usepackage[utf8]{inputenc}
\usepackage{lmodern}
\usepackage{microtype}
\usepackage{parskip}
\usepackage{enumitem}
\usepackage{booktabs}
\usepackage{tabularx}
\usepackage{longtable}
\usepackage{array}
\usepackage{xcolor}
\usepackage[most]{tcolorbox}
\usepackage{fancyhdr}
\usepackage{hyperref}
\usepackage{amsmath,amssymb}
\usepackage{listings}
\usepackage{tikz}
\usetikzlibrary{arrows.meta,positioning,shapes.geometric,fit}
\usepackage{needspace}
\usepackage{multicol}
\usepackage{lastpage}
\usepackage{titlesec}

\definecolor{navy}{HTML}{17324D}
\definecolor{bluegray}{HTML}{43657D}
\definecolor{lightblue}{HTML}{EAF2F8}
\definecolor{lightgray}{HTML}{F3F5F7}
\definecolor{linegray}{HTML}{C9D1D9}
\definecolor{darkgray}{HTML}{4A5561}
\definecolor{accent}{HTML}{8E3B46}

\hypersetup{colorlinks=true,linkcolor=navy,urlcolor=bluegray,citecolor=navy,pdfborder={0 0 0}}
\setlist[itemize]{leftmargin=1.4em,itemsep=2pt,topsep=4pt}
\setlist[enumerate]{leftmargin=1.6em,itemsep=2pt,topsep=4pt}
\renewcommand{\arraystretch}{1.22}
\setlength{\tabcolsep}{4pt}
\newcolumntype{Y}{>{\raggedright\arraybackslash}X}
\newcolumntype{C}[1]{>{\centering\arraybackslash}p{#1}}
\newcolumntype{L}[1]{>{\raggedright\arraybackslash}p{#1}}
\pagestyle{fancy}
\fancyhf{}
\lhead{Application Security Engineering}
\rhead{Code Review Handbook}
\cfoot{\thepage\ of \pageref{LastPage}}
\renewcommand{\headrulewidth}{0.4pt}
\titleformat{\section}{\Large\bfseries\color{navy}}{\thesection}{0.65em}{}
\titleformat{\subsection}{\large\bfseries\color{bluegray}}{\thesubsection}{0.55em}{}
\titleformat{\subsubsection}{\normalsize\bfseries\color{darkgray}}{\thesubsubsection}{0.5em}{}
\titlespacing*{\section}{0pt}{1.4em}{0.55em}
\titlespacing*{\subsection}{0pt}{1.0em}{0.35em}

\newcommand{\checkbox}{\ensuremath{\square}}
\newcommand{\fieldline}[1]{\textbf{#1}\hspace{0.5em}\leaders\hbox{\rule[0.4ex]{0.6em}{0.35pt}\hspace{0.25em}}\hfill\kern0pt}
\newcommand{\workarea}[1]{\par\smallskip\noindent\begin{tcolorbox}[height=#1,colback=white,colframe=linegray,boxrule=0.5pt,arc=0pt,left=2pt,right=2pt,top=2pt,bottom=2pt]\mbox{}\end{tcolorbox}\par}
\newtcolorbox{keybox}[1]{colback=lightblue,colframe=bluegray,title=\textbf{#1},fonttitle=\bfseries,boxrule=0.7pt,arc=2pt,left=6pt,right=6pt,top=5pt,bottom=5pt}
\newtcolorbox{cautionbox}[1]{colback=lightgray,colframe=accent,title=\textbf{#1},fonttitle=\bfseries,boxrule=0.7pt,arc=2pt,left=6pt,right=6pt,top=5pt,bottom=5pt}
\newtcolorbox{worksheetbox}[1]{colback=white,colframe=linegray,title=\textbf{#1},fonttitle=\bfseries,coltitle=navy,boxrule=0.6pt,arc=1pt,left=6pt,right=6pt,top=5pt,bottom=5pt}
\lstset{basicstyle=\ttfamily\small,columns=fullflexible,breaklines=true,frame=single,framerule=0.4pt,rulecolor=\color{linegray},backgroundcolor=\color{lightgray},showstringspaces=false,tabsize=2,keywordstyle=\bfseries}
\setcounter{tocdepth}{2}

\title{\textbf{Desk-Check for Algorithms Workbook}\\\large Application Security Engineering Handout}
\author{}
\date{}

\begin{document}
\maketitle
\thispagestyle{empty}
\begin{keybox}{Purpose}
Hands-on workbook for manual algorithm execution, boundary analysis, variable relationships, and security-state validation.
\end{keybox}
\vfill
\noindent\textbf{Use note.} This handout is a practitioner aid for authorized software security assessment. It emphasizes evidence, traceability, and disciplined review rather than a single rigid methodology.
\newpage
\tableofcontents
\newpage

\section{Desk-Check Method}
Desk-checking is manual execution on paper (or in a worksheet). The reviewer chooses suspicious initial values, records relevant state, and steps through the algorithm exactly as the implementation would. The method is especially useful when variable relationships, repeated assignment, loop termination, side effects, or boundary transitions are difficult to simulate reliably in one's head.

\begin{keybox}{Core Procedure}
\begin{enumerate}
\item Select a bounded code fragment or algorithm and state its intended purpose.
\item Inventory inputs, outputs, mutable state, external state, invariants, and security expectations.
\item Choose high-value initial values, emphasizing boundaries and unusual state combinations.
\item Execute statement by statement. Record branches, mutations, allocations, resource changes, and side effects.
\item Inspect the state immediately after loops, error paths, and termination points.
\item Compare the resulting state with the expected postconditions and security invariants.
\item Record a candidate finding or a new test hypothesis when observed behavior is unsafe or unexplained.
\end{enumerate}
\end{keybox}

\section{Pre-Desk-Check Analysis}
\begin{worksheetbox}{Algorithm Profile}
\par\smallskip\noindent\begin{tabularx}{\linewidth}{@{}L{0.29\linewidth}Y@{}}
Algorithm / function & \\ \midrule
File / lines & \\ \midrule
Purpose & \\ \midrule
Inputs & \\ \midrule
Outputs & \\ \midrule
Mutable local state & \\ \midrule
External / shared state & \\ \midrule
Preconditions & \\ \midrule
Postconditions & \\ \midrule
Security invariants & \\ \midrule
Security expectations & \\ \midrule
Attacker-controlled variables & \\ \midrule
Trusted / constrained variables & \\ \midrule
Derived variables & \\ \midrule
Buffers / objects / resources & \\ 
\end{tabularx}
\end{worksheetbox}
\clearpage

\section{Test Case Selection}
Boundary cases provide high information density, but the most important boundaries are those created by the implementation itself. A value can become dangerous because of arithmetic, rounding, normalization, truncation, conversion, indexing, or interaction with another variable even when it is not an obvious domain boundary.

\begin{longtable}{@{}L{0.27\linewidth}L{0.25\linewidth}L{0.18\linewidth}L{0.23\linewidth}@{}}
\toprule
Case family & Candidate value / state & Applicable? & Why it matters here \\
\midrule
Zero / empty & & & \\
Minimum representable / valid & & & \\
Maximum representable / valid & & & \\
Just below a boundary & & & \\
Exactly at a boundary & & & \\
Just above a boundary & & & \\
Small negative & & & \\
Large negative & & & \\
Small positive & & & \\
Large positive & & & \\
Null / absent & & & \\
Malformed & & & \\
Truncated & & & \\
Oversized & & & \\
Repeated / duplicate & & & \\
Unusual encoding / representation & & & \\
Unexpected state combination & & & \\
Implementation-created boundary & & & \\
\bottomrule
\end{longtable}
\clearpage

\section{Primary Desk-Check Execution Table}
Use as many variable-state entries as necessary; do not force the algorithm into a fixed number of columns.
\small
\begin{longtable}{@{}C{0.06\linewidth}L{0.22\linewidth}L{0.14\linewidth}L{0.20\linewidth}L{0.18\linewidth}L{0.12\linewidth}@{}}
\toprule
Step & Statement / event & Branch result & Variable state & Memory / object state & Security observation \\
\midrule
\endfirsthead
\toprule
Step & Statement / event & Branch result & Variable state & Memory / object state & Security observation \\
\midrule
\endhead
\rule{0pt}{2.7em} & & & & & \\
\rule{0pt}{2.7em} & & & & & \\
\rule{0pt}{2.7em} & & & & & \\
\rule{0pt}{2.7em} & & & & & \\
\rule{0pt}{2.7em} & & & & & \\
\rule{0pt}{2.7em} & & & & & \\
\rule{0pt}{2.7em} & & & & & \\
\rule{0pt}{2.7em} & & & & & \\
\bottomrule
\end{longtable}
\normalsize

\subsection{Compact Repetition Sheet}
\par\smallskip\noindent\begin{tabularx}{\linewidth}{@{}C{0.07\linewidth}L{0.27\linewidth}L{0.19\linewidth}Y@{}}
\toprule
Step & Statement / condition & State delta & Observation \\
\midrule
\rule{0pt}{2.5em} & & & \\
\rule{0pt}{2.5em} & & & \\
\rule{0pt}{2.5em} & & & \\
\rule{0pt}{2.5em} & & & \\
\rule{0pt}{2.5em} & & & \\
\bottomrule
\end{tabularx}

\section{Loop Analysis Worksheet}
\par\smallskip\noindent\begin{tabularx}{\linewidth}{@{}L{0.29\linewidth}Y@{}}
Initialization & \\ \midrule
Loop condition & \\ \midrule
Loop variable(s) & \\ \midrule
Increment / decrement & \\ \midrule
Expected termination condition & \\ \midrule
Maximum expected iteration & \\ \midrule
Boundary transition & \\ \midrule
Index value(s) & \\ \midrule
Target object / capacity & \\ \midrule
Off-by-one risk & \\ \midrule
State on final body iteration & \\ \midrule
\textbf{State immediately after termination} & \\ \midrule
Cleanup / termination operation & \\ 
\end{tabularx}

\begin{cautionbox}{Post-Loop Check}
A loop can preserve its intended invariant during every body iteration and still leave a dangerous index, length, pointer, or state value for the statement that runs immediately after the loop. Always record post-loop state explicitly.
\end{cautionbox}

\section{Branch Analysis}
\small
\begin{longtable}{@{}L{0.20\linewidth}L{0.17\linewidth}L{0.17\linewidth}C{0.09\linewidth}L{0.14\linewidth}L{0.15\linewidth}@{}}
\toprule
Condition & True-path state & False-path state & Reachable? & Attacker influence & Exceptional combination / consequence \\
\midrule
\endfirsthead
\toprule
Condition & True-path state & False-path state & Reachable? & Attacker influence & Exceptional combination / consequence \\
\midrule
\endhead
\rule{0pt}{2.7em} & & & & & \\
\rule{0pt}{2.7em} & & & & & \\
\rule{0pt}{2.7em} & & & & & \\
\rule{0pt}{2.7em} & & & & & \\
\bottomrule
\end{longtable}
\normalsize
\clearpage

\section{Memory and Resource State}
Track only fields relevant to the algorithm.
\par\smallskip\noindent\begin{tabularx}{\linewidth}{@{}L{0.22\linewidth}L{0.18\linewidth}L{0.18\linewidth}Y@{}}
\toprule
State item & Initial & Current / final & Security question \\
\midrule
Buffer size / capacity & & & Is every write bounded by actual capacity? \\
Index / offset & & & Can it cross the valid interval? \\
Allocation size & & & Can arithmetic wrap, truncate, or under-allocate? \\
Bytes used / remaining & & & Do relationships remain consistent? \\
Ownership & & & Who may free / close / mutate? \\
Lifetime & & & Can use occur before init or after release? \\
Resource state & & & Are failure and cleanup paths balanced? \\
Lock state & & & Are acquisition / release assumptions preserved? \\
Privilege state & & & Can state cross a trust boundary unexpectedly? \\
\bottomrule
\end{tabularx}

\section{Variable Relationship Worksheet}
Many failures arise not from one variable's value but from an invalid relationship between two or more values.
\begin{longtable}{@{}L{0.14\linewidth}L{0.14\linewidth}L{0.23\linewidth}L{0.18\linewidth}L{0.24\linewidth}@{}}
\toprule
Variable A & Variable B & Expected relationship & Enforced where? & Attacker influence / boundary / consequence \\
\midrule
\endfirsthead
\toprule
Variable A & Variable B & Expected relationship & Enforced where? & Attacker influence / boundary / consequence \\
\midrule
\endhead
length & capacity & $0 \leq length \leq capacity$ & & \\
index & upper bound & index remains within the valid element range & & \\
count & element size & multiplication and allocation remain representable & & \\
declared length & actual data & parser never trusts one without reconciling the other & & \\
state flag & object lifecycle & state matches initialization / ownership reality & & \\
\rule{0pt}{2.5em} & & & & \\
\rule{0pt}{2.5em} & & & & \\
\bottomrule
\end{longtable}
\clearpage

\section{Security Invariant Register}
\begin{longtable}{@{}L{0.23\linewidth}L{0.16\linewidth}L{0.16\linewidth}L{0.18\linewidth}L{0.20\linewidth}@{}}
\toprule
Invariant & Established at & Consumed at & How it can fail & Validation / result \\
\midrule
\endfirsthead
\toprule
Invariant & Established at & Consumed at & How it can fail & Validation / result \\
\midrule
\endhead
\rule{0pt}{2.6em} & & & & \\
\rule{0pt}{2.6em} & & & & \\
\rule{0pt}{2.6em} & & & & \\
\rule{0pt}{2.6em} & & & & \\
\bottomrule
\end{longtable}

\section{Constraint Establishment and Test-Space Reduction}
Discard a state only after demonstrating that it is unreachable under the reviewed implementation. Validate the enforcing check separately; future or alternate callers can invalidate an assumption that is true for one path today.

\par\smallskip\noindent\begin{tabularx}{\linewidth}{@{}L{0.20\linewidth}L{0.20\linewidth}L{0.18\linewidth}L{0.18\linewidth}Y@{}}
\toprule
Candidate state & Claimed constraint & Enforcement location & Bypass / alternate caller? & Keep or discard + evidence \\
\midrule
\rule{0pt}{2.7em} & & & & \\
\rule{0pt}{2.7em} & & & & \\
\rule{0pt}{2.7em} & & & & \\
\bottomrule
\end{tabularx}

\subsection{Extraneous Input Thinning}
An input state may be deprioritized only when the reviewer has evidence that (1) an attacker cannot influence it directly or indirectly, and (2) varying it does not materially alter processing of attacker-controlled data or security-sensitive behavior.
\par\smallskip\noindent\begin{tabularx}{\linewidth}{@{}L{0.16\linewidth}L{0.17\linewidth}C{0.12\linewidth}L{0.18\linewidth}Y@{}}
\toprule
Variable & Source & Attacker control? & Behavioral impact & Keep / remove + rationale \\
\midrule
\rule{0pt}{2.5em} & & & & \\
\rule{0pt}{2.5em} & & & & \\
\rule{0pt}{2.5em} & & & & \\
\bottomrule
\end{tabularx}
\clearpage

\section{Unconstrained Data Lifecycle}
For strings, structured text, serialized objects, and other broad input domains, classify each operation by what it does to the representation.
\begin{longtable}{@{}L{0.18\linewidth}L{0.20\linewidth}L{0.24\linewidth}L{0.31\linewidth}@{}}
\toprule
Stage & Operation & Security questions & Test ideas \\
\midrule
Transformation & Decode, normalize, recode, truncate, expand & Can representation change invalidate an earlier assumption? Does ordering relative to validation matter? & Alternate encodings, normalization collisions, truncation boundaries, expansion limits \\
Validation & Syntax, length, allow/deny rules, semantic checks & Can the check be bypassed? Is it sufficient for later use? & Boundary lengths, mixed representations, incomplete checks \\
Parsing / extraction & Delimit, split, deserialize, identify fields & Do parser states agree on boundaries and lengths? & Empty fields, nested delimiters, inconsistent lengths, malformed structure \\
System / resource use & Select file, query, command, identity, resource & Can the value select an unintended resource or cross a trust boundary? & Alternate resource names, unexpected namespace, privilege-sensitive target \\
\bottomrule
\end{longtable}

\section{Desk-Check Finding Capture}
\begin{worksheetbox}{Observation Record}
\par\smallskip\noindent\begin{tabularx}{\linewidth}{@{}L{0.29\linewidth}Y@{}}
Observation & \\ \midrule
Violated assumption / invariant & \\ \midrule
Code location & \\ \midrule
Triggering values & \\ \midrule
Resulting state & \\ \midrule
Security consequence & \\ \midrule
Confidence & \\ \midrule
Next validation step & \\ 
\end{tabularx}
\end{worksheetbox}
\clearpage

\section{Practice Exercises}
The examples below are original and intentionally small. Perform the worksheet before reading the answer key.

\subsection{Exercise 1 -- Off-by-One Write}
\begin{lstlisting}[language=C]
int copy_token(char *dst, size_t cap, const char *src) {
    size_t i;
    for (i = 0; src[i] != '\0' && i <= cap; ++i) {
        dst[i] = src[i];
    }
    dst[i] = '\0';
    return 0;
}
\end{lstlisting}
\textbf{Questions.} What happens when the source length equals \texttt{cap}? What is the post-loop value of \texttt{i}? Which write is outside the intended destination region?
\workarea{0.9in}

\subsection{Exercise 2 -- Integer Boundary in Allocation}
\begin{lstlisting}[language=C]
int reserve_items(size_t count, size_t item_size) {
    size_t bytes = count * item_size;
    void *p = malloc(bytes + 16);
    if (!p) return -1;
    /* use p for count items */
    free(p);
    return 0;
}
\end{lstlisting}
\textbf{Questions.} Which variable relationship matters? What values near the multiplication and addition boundaries should be tested? Can a small resulting allocation represent a much larger logical request?
\workarea{0.9in}

\subsection{Exercise 3 -- Validation Before Transformation}
\begin{lstlisting}[language=Python]
def open_report(name):
    if "../" in name:
        return None
    normalized = urllib.parse.unquote(name)
    return open("/srv/reports/" + normalized, "rb")
\end{lstlisting}
\textbf{Questions.} Which operation is validation and which is transformation? Can the transformation create a sequence that validation never observed? What ordering would preserve the intended invariant?
\workarea{0.9in}

\subsection{Exercise 4 -- Unchecked State Transition}
\begin{lstlisting}[language=C]
enum state { NEW, AUTHENTICATED, CLOSED };

int perform(struct session *s, int privileged_action) {
    if (privileged_action && s->state != CLOSED) {
        return run_admin_action(s);
    }
    return 0;
}
\end{lstlisting}
\textbf{Questions.} Which states permit the privileged action? Which state was probably intended? Which security invariant is missing from the condition?
\workarea{0.9in}

\subsection{Exercise 5 -- Buffer/Index Relationship}
\begin{lstlisting}[language=C]
int read_byte(struct buf *b, size_t offset) {
    if (offset > b->used) return -1;
    return b->data[offset];
}
\end{lstlisting}
\textbf{Questions.} What does \texttt{used} represent? If it is a count of valid bytes, what is the valid maximum index? Which boundary value should be desk-checked?
\workarea{0.9in}

\subsection{Exercise 6 -- State-Dependent Error Path}
\begin{lstlisting}[language=C]
int update(struct obj *o, const char *v) {
    lock(o->mu);
    o->busy = 1;
    if (!validate(v))
        return -1;
    apply(o, v);
    o->busy = 0;
    unlock(o->mu);
    return 0;
}
\end{lstlisting}
\textbf{Questions.} What state remains after validation fails? What resource state remains? What later behavior could depend on \texttt{busy} or the lock?
\workarea{0.9in}
\clearpage

\section{Answer Key}
\subsection*{Exercise 1}
The loop admits \texttt{i == cap}, so the body can write \texttt{dst[cap]}. After that iteration \texttt{i} becomes \texttt{cap + 1}; the terminator write can be even farther outside the intended region. The exact outcome depends on how capacity is defined, which is why the relationship between capacity and valid index range must be stated explicitly.

\subsection*{Exercise 2}
The critical relationship is \texttt{count * item\_size} plus the fixed overhead. Both arithmetic operations require representability checks before allocation. Values around \(\lfloor SIZE\_MAX/item\_size\rfloor\) and around the final addition boundary are high-value cases.

\subsection*{Exercise 3}
The substring check validates the pre-decoded representation; URL decoding is a later transformation. Encoded traversal markers can appear only after decoding. Validate the canonical representation that will actually be used, and ensure the resource resolution itself cannot escape the intended namespace.

\subsection*{Exercise 4}
The condition permits both \texttt{NEW} and \texttt{AUTHENTICATED}. If privileged actions require authentication, the invariant should require the authenticated state explicitly rather than merely excluding \texttt{CLOSED}.

\subsection*{Exercise 5}
If \texttt{used} is a count, valid indices are \(0\) through \(used-1\). The condition should reject \texttt{offset == used}; desk-check that exact boundary.

\subsection*{Exercise 6}
The error path returns while the object remains busy and the lock remains held. This can create denial of service, deadlock, or inconsistent lifecycle behavior. The desk-check should track both logical state and resource state across every early return.
\clearpage

\section{Reusable Blank Desk-Check Sheet}
\begin{worksheetbox}{Case Metadata}
\fieldline{Function / algorithm:}\par
\fieldline{Case ID:}\par
\fieldline{Initial values / assumptions:}\par
\fieldline{Expected invariant / outcome:}
\end{worksheetbox}
\begin{longtable}{@{}C{0.06\linewidth}L{0.24\linewidth}L{0.24\linewidth}L{0.20\linewidth}L{0.19\linewidth}@{}}
\toprule
Step & Statement / branch & State snapshot & Resource / object state & Observation \\
\midrule
\endfirsthead
\toprule
Step & Statement / branch & State snapshot & Resource / object state & Observation \\
\midrule
\endhead
\rule{0pt}{2.8em} & & & & \\
\rule{0pt}{2.8em} & & & & \\
\rule{0pt}{2.8em} & & & & \\
\rule{0pt}{2.8em} & & & & \\
\rule{0pt}{2.8em} & & & & \\
\rule{0pt}{2.8em} & & & & \\
\rule{0pt}{2.8em} & & & & \\
\rule{0pt}{2.8em} & & & & \\
\bottomrule
\end{longtable}
\end{document}
